% PARADIGM.md section 2 (Theorem A).  File id: thmA.
\section{Theorem A: identification}\label{sec:thmA}

\begin{thm}{Theorem}{A}
Assume \textup{(S0)}.
\begin{enumerate}[label=(\roman*)]
  \item Assume scalar classical gains $\Gamma_W = \gamma_W I$, $\Gamma_c = \gamma_c I$ and
    $\Gamma_\sigma = \gamma_\sigma I$, and use the dictionary $W_i = a_i/m$, $\gamma_W = g_a/m$, $\gamma_c = m g_c$,
    $\gamma_\sigma = m g_s$.
    \begin{itemize}
      \item The \emph{raw} \textup{(}unprojected\textup{)} rates of \textup{(L)} and of the classical law of
        Section~\ref{sec:p1-algorithm} coincide.
      \item The \emph{projected} rates coincide when the classical joint projections are replaced by per-kernel
        ones. The classical law projects $W$ onto one Frobenius ball and $\operatorname{vec}(c)$ onto one Euclidean
        ball. The per-kernel version projects:
        \begin{itemize}
          \item $W_i$ onto $\norm{W_i}\le a_{\max}/m$, with layer $\varepsilon_K/m$;
          \item $c_i$ onto $\norm{c_i}\le c_{\max}$, with layer $\varepsilon_K$;
          \item $\sigma_i$ onto $[s_{\min},s_{\max}]$, with layer $\varepsilon_K$.
        \end{itemize}
      \item For $m\ge 2$, no joint ball, of any radius $W_{\max}$ and layer $\varepsilon_W$, reproduces the
        per-kernel field. For $m = 1$ the two coincide, with $W_{\max} = a_{\max}$ and layer $\varepsilon_K$.
    \end{itemize}
  \item Fix a closed-loop trajectory $(x(\cdot),e(\cdot))$. Then:
    \begin{itemize}
      \item The field $v_t(\vartheta) := \Proj(\vartheta, G\nabla g_t(\vartheta))$ is Lipschitz in $\vartheta$ on
        $\ThK$, and $\ThK$ is forward invariant.
      \item For every $\mu_0\in\Prob(\ThK)$, $\mu_t := (X_t)_\#\mu_0$ is a distributional solution of
        $\partial_t\mu + \nabla\cdot(\mu v_t) = 0$.
      \item On $\operatorname{int}\ThKe$, $v_t = -G\nabla(\delta F_t/\delta\mu)$. So at frozen time, $\mu_t$ is the
        Wasserstein gradient flow \textup{(}metric $\Ginv$\textup{)} of the linear functional $F_t$, and
        \textup{(L)} is its particle form.
      \item $F_t$ is generated by the closed loop, so the coupled system is the one studied in Theorem D.
    \end{itemize}
  \item \textbf{Dictionary}, all in the $\Ginv$ metric:

    \smallskip
    {\normalfont\small
    \begin{tabularx}{\linewidth}{@{}>{\hsize=0.8\hsize\raggedright\arraybackslash}X
                                    >{\hsize=1.2\hsize\raggedright\arraybackslash}X@{}}
      \toprule
      classical device & measure-level object\\
      \midrule
      projection onto $\ThK$ &
      A velocity-level (Pomet--Praly) regularisation of the support constraint $\mu_t\in\Prob(\ThK)$. For
      $\varepsilon_K>0$ it is not the gradient flow of any functional of $\mu$. Formally, as $\varepsilon_K\to 0$,
      the field tends to the $\Ginv$-projection of $G\nabla g_t$ onto the tangent cone of $\ThK$, which is the
      velocity of the $\WG$-gradient flow of $F_t + \iota_{\Prob(\ThK)}$ (\cite[Prop.~3.1]{R46} for the exact $W_2$
      projection).\\
      \addlinespace
      $\sigma$-modification in the $\Ginv$ metric (leak outside the gain,
      $\dot\vartheta = G\nabla g - \sigma(\vartheta-\vartheta^0)$, as in $\Lsig$ and Theorem H) &
      the potential energy $(\sigma/2)\int\norm{\vartheta-\vartheta^0}_{\Ginv}^2\dd\mu$, which is geodesically
      $\sigma$-convex\\
      \addlinespace
      classical $\sigma$-modification $\Gamma(y-\sigma\theta)$ (leak inside the gain;
      \code{src/baselines/sigma_mod.py} implements its $W$-block, with $c$ and $\sigma$ frozen) &
      For the full-vector leak, under (i), the object is
      $(\sigma/2)\int(\vartheta-\vartheta^0)^\T D(\vartheta-\vartheta^0)\dd\mu$ with
      $D = \operatorname{diag}(I_n/m, mI_n, m)$. In the $\Ginv$ metric it is $\lambda$-convex with
      $\lambda = \sigma\cdot\min(g_a/m, mg_c, mg_s)$, which depends on $m$. For the $W$-block alone it is
      $(\sigma/2)\int(a-a^0)^\T(I_n/m)(a-a^0)\dd\mu$.\\
      \addlinespace
      $\sigma$-modification toward a labelled $(\vartheta^0_i)$ &
      not a functional of $\mu$; it is a functional on the labelled lift $\R^{mp}$\\
      \addlinespace
      OT $\sigma$-modification (Theorem C) &
      $(\sigma/2)W(\mu,\mu^0)^2$. For atomic $\mu$ with a unique optimal assignment, its $W$-gradient at atom $i$ is
      $\sigma(\vartheta_i-\vartheta^0_{\tau(i)})$. At ties, $\Lsig$ uses a measurable selection from the
      superdifferential $\{\sigma(\vartheta_i-\vartheta^0_{\tau(i)}) : \tau\text{ optimal}\}$.\\
      \addlinespace
      $e$-modification & the same functionals with coefficient $\sigma\norm{e}$\\
      \addlinespace
      parameter dither (Theorem H) & the entropy $\tau\int\mu\log\mu$\\
      \bottomrule
    \end{tabularx}}
\end{enumerate}
\end{thm}
\status{Proved and checked. Novelty (literature audit): interpretive; ingredients known
\cite{R29}, \cite{R41}, \cite{R46}, \cite{R25}.}
\checkedby{\code{tests/test_transport.py}: \code{test_identification_with_classical_law} (raw rates in (i));\newline
\code{test_identification_projected_rates_per_kernel} (projected rates and the only-if construction in (i)).}

\begin{proof}
\textbf{(i)} \textbf{Gain map.} We have $\nabla_a g = \varphi e$, $\nabla_c g = (a^\T e)\nabla_c\varphi$ and
$\partial_s g = (a^\T e)\partial_s\varphi$. Matching these with the classical rates gives the gain map.

\textbf{Scaling identity.}
$(1/m)\operatorname{Proj}_{r,\varepsilon}(a,y) = \operatorname{Proj}_{r/m,\varepsilon/m}(a/m, y/m)$. Three facts give
it:
\begin{itemize}
  \item $h_{r/m,\varepsilon/m}(\norm{a}/m) = h_{r,\varepsilon}(\norm{a})$;
  \item the direction $a(a^\T y)/\norm{a}^2$ is homogeneous of degree $0$ in $a$ and of degree $1$ in $y$;
  \item the activation $a^\T y>0$ is invariant under scaling.
\end{itemize}
If the layer is not rescaled, the ramps differ.

\emph{Only if \textup{(}$m\ge 2$\textup{)}.} Let $r_{\mathrm{in}} := (a_{\max}-\varepsilon_K)/m$, and let the joint
ball have radius $W_{\max}$, layer $\varepsilon_W$ and $\rho := W_{\max}-\varepsilon_W$. In both cases below the
driving term points outward.
\begin{itemize}
  \item If $\rho<\sqrt2\,r_{\mathrm{in}}$, take two rows of equal norm $r$ with
    $\sqrt2\,r\in\bigl(\rho,\min(\sqrt2\,r_{\mathrm{in}},W_{\max})\bigr]$, parallel to $e$. The per-kernel field is
    inactive and the joint field is active.
  \item If $\rho\ge\sqrt2\,r_{\mathrm{in}}$, take one row with norm in
    $\bigl(r_{\mathrm{in}},\min(a_{\max}/m,\rho)\bigr)$ and the others zero. The per-kernel field is active and the
    joint field is inactive.
\end{itemize}
The same argument applies to $\operatorname{vec}(c)$. For $m = 1$ the two operators coincide.

The test \code{test_identification_projected_rates_per_kernel} checks three things: the identity with active projections; that the layer must be
rescaled; and the only-if construction against \code{ProjectionOperator.project_matrix} on five joint balls. Four are
active against an inactive per-kernel field, and one is inactive against an active per-kernel field.

\textbf{(ii)} $\Proj$ is Lipschitz in $\vartheta$. The ramp vanishes for
$\norm{\vartheta}\le\theta_{\max}-\varepsilon_K$, so the direction $\vartheta/\norm{\vartheta}$ is evaluated only away
from $0$. Also, $\nabla g_t$ is $C^1$ on $s\ge s_{\min}$.
\begin{itemize}
  \item \textbf{Invariance.} Forward invariance follows from Nagumo's theorem, since
    $n_b^\T\operatorname{Proj}_b\le 0$ on the boundary, with the outward normals $n_b$ of Lemma P.1.
  \item \textbf{Continuity equation.} The push-forward of $\mu_0$ by a Lipschitz flow solves the continuity equation
    (Golse~\cite{R82}; Ambrosio--Gigli--Savaré~\cite[\S8.1]{ags2008}).
  \item \textbf{Gradient form.} On $\operatorname{int}\ThKe$ the projection is inactive, so
    $v_t = G\nabla g_t = -G\nabla(\delta F_t/\delta\mu)$.
\end{itemize}

\textbf{(iii)} \textbf{Potential-energy rows.} These follow from the computation in the metric $\Ginv$: for any $V$,
$\nabla_{\Ginv}V = G\nabla V$. For example, $V = (\sigma/2)\norm{\vartheta-\vartheta^0}_{\Ginv}^2$ gives
$G\nabla V = \sigma(\vartheta-\vartheta^0)$.

\textbf{Classical-leak row.} Applying the gain map of (i) to the classical leak $\Gamma(y-\sigma\theta)$ gives the
velocity $-\sigma GD(\vartheta-\vartheta^0)$.

\textbf{Projection row.} The $\Ginv$-rescaled correction field has an asymmetric Jacobian. Take the example $m = 1$,
$n = 1$, $G = I$, $x = 0.7$, $e = 1$, $\vartheta = (0.93, 0.2, 1.5)$, $a_{\max} = 1$, $\varepsilon_K = 0.15$ (so
$h(0.93) = 0.284$) and $c_{\max} = 2$, with $a$ in the layer. A numerical evaluation gives:
\begin{itemize}
  \item the $c$-derivative of the $a$-component is $-0.060$;
  \item the $a$-derivative of the $c$-component is $0$.
\end{itemize}
So the field is not a gradient.
\end{proof}

\begin{rmk}{Remark}{A.1}[normalisation]
The algorithm's raw rates are unchanged. Two things do change: the admissible set becomes a product of per-kernel
balls, and $a_{\max}$, not $\norm{W}_F$, is held fixed in $m$.
\begin{itemize}
  \item \textbf{$\bm{m}$-free constants.} In this normalisation the structural constants $k_e$, $K$, $L_0$, $L_1$,
    $\lambda_G$ and $\Delta$ are free of $m$ for fixed $G$ and $a_{\max}$. So are the constants of Theorems B,
    C(i)--(iii), D and F.
  \item \textbf{$\bm{m}$-dependent data.} $\varepsilon_{\max}$ and $V_*$ are data of the chosen $m$-atom ideal. The
    local constants of C(iv) and of Theorem E depend on $m$.
  \item \textbf{Classical convention.} Fixed classical $\gamma$'s correspond to $g_a = m\gamma_W\propto m$ and
    $g_c, g_s\propto 1/m$. Weights speed up, centres and widths freeze, $\lambda_G\propto m$ and
    $\Delta\propto\sqrt m$.
\end{itemize}
\end{rmk}
